\BourbakiExercises{习题}

\begin{enumerate}
\def\labelenumi{\arabic{enumi})}
\tightlist
\item
  设 A、B 是环，P 是一个 (A, B)-双模，Q 是一个 (B, A)-双模。设 \(\varphi: P \otimes_{B} Q \to A\) 是一个 (A, A)-双模同态，\(\psi: Q \otimes_{A} P \to B\) 是一个 (B, B)-双模同态，满足对 P 中的 x, y 与 Q 中的 u, v 有 \(\varphi(x \otimes u)y = x\psi(u \otimes y)\) 及 \(u\varphi(x \otimes v) = \psi(u \otimes x)v\)。设 \(\varphi\) 是满射。
\end{enumerate}

\begin{enumerate}
\def\labelenumi{\alph{enumi})}
\item
  证明 A-模 P 是生成模。
\item
  证明 \(\varphi\) 是同构。
\item
  构造 Q 与 \(\operatorname{Hom}_{\mathbb{B}}(\mathrm{P}, \mathrm{B})\) 之间的一个 (B, A)-双模同构。
\end{enumerate}

\begin{enumerate}
\def\labelenumi{\arabic{enumi})}
\setcounter{enumi}{1}
\item
  设 M、N 是 A-模，m、n 是 \(\geqslant 1\) 的整数。设 M 是 \(N^{n}\) 的一个直因子，N 是 \(M^{m}\) 的一个直因子；证明环 \(\operatorname{End}_{\mathrm{A}}(\mathrm{M})\) 与 \(\operatorname{End}_{\mathrm{A}}(\mathrm{N})\) 是森田等价的。
\item
  设 A 是一个环，G 是一个用加法记号表示的交换群。映射 \(T: A \to G\) 称为迹型的，若它是一个加群同态并且对 A 中的所有 a, b 有 \(T(ab) = T(ba)\)。
\end{enumerate}

\begin{enumerate}
\def\labelenumi{\alph{enumi})}
\item
  设 \(T : A \to G\) 是一个迹型映射，P 是一个投射的且有限生成的右 A-模。证明存在唯一的迹型映射 \(\mathrm{T}_{\mathrm{P}} : \mathrm{End}_{\mathrm{A}}(\mathrm{P} \oplus \mathrm{A}_{d}) \to \mathrm{G}\)，它在 \(\mathrm{A} = \mathrm{End}_{\mathrm{A}}(\mathrm{A}_{d})\) 上的限制为 T。
\item
  设 A 与 B 是森田等价的环，G 是一个交换群。构造从 A 到 G 的迹型映射与从 B 到 G 的迹型映射之间的一个双射对应。
\end{enumerate}

\begin{enumerate}
\def\labelenumi{\arabic{enumi})}
\setcounter{enumi}{3}
\tightlist
\item
  设 A 与 B 是环，P 是一个可逆 (A,B)-双模。用 \(\mathcal{B}_{\mathrm{A}}\)（相应地，\(B_{B}\)）表示 A（相应地，B）的双边理想之集，用 \(f: B_{A} \to B_{B}\) 表示命题 6（VIII, 第 105 页）中定义的有序集同构。
\end{enumerate}

\begin{enumerate}
\def\labelenumi{\alph{enumi})}
\item
  设 \(\mathfrak{a}\) 是 A 的一个理想。证明环 A/ \(\mathfrak{a}\) 与 B/f( \(\mathfrak{a}\) ) 是森田等价的。
\item
  设 \(\mathfrak{a}\) 与 \(\mathfrak{a}'\) 是 A 的理想。证明 \(f(\mathfrak{aa}') = f(\mathfrak{a})f(\mathfrak{a}')\)。
\item
  证明 A 的理想 \(\mathfrak{a}\) 是幂零的当且仅当 B 的理想 \(f(\mathfrak{a})\) 是幂零的。 d) 设 M 是一个 B-模。证明 M 的迹理想（VIII, 第 80 页）是在 \(f\) 之下 \(\mathrm{P} \otimes_{\mathrm{B}} \mathrm{M}\) 的迹理想的像。
\end{enumerate}

¶ 5) 对一个环 R，用 \(\mathbf{M}_{\infty}(\mathrm{R})\) 表示由 \(\mathbf{N} \times \mathbf{N}\) 型、系数在 R 中且只有有限个非零系数的矩阵构成的伪环。证明环 A 与 B 森田等价当且仅当伪环 \(\mathbf{M}_{\infty}(\mathrm{A})\) 与 \(\mathbf{M}_{\infty}(\mathrm{B})\) 同构。（对一个可逆 (A,B)-模 P，考虑 \(\mathbf{M}_{\infty}(\mathrm{A})\) 在 VIII, 第 92 页习题 12 中构造的同构 \(\mathrm{End}_{\mathrm{A}}(\mathrm{A}_{s}^{(\mathbf{N})}) \to \mathrm{End}_{\mathrm{A}}(\mathrm{P}^{(\mathbf{N})})\) 下的像。给定一个从 \(\mathbf{M}_{\infty}(\mathrm{A})\) 到 \(\mathbf{M}_{\infty}(\mathrm{B})\) 的同构，考虑在 \(\mathbf{M}_{\infty}(\mathrm{B})\) 中幂等元 \(\mathrm{E}_{11}\)（属于 \(\mathbf{M}_{\infty}(\mathrm{A})\)）的像。）

\begin{enumerate}
\def\labelenumi{\arabic{enumi})}
\setcounter{enumi}{5}
\tightlist
\item
  设 \(\mathrm{A}\) 与 \(\mathrm{B}\) 是环。设给定了
\end{enumerate}

- 对每个 A-模 M，一个 B-模 \(\mathrm{F(M)}\)，以及

- 一个 B-线性同态 \(\mathrm{F}(u): \mathrm{F}(\mathrm{M}) \to \mathrm{F}(\mathrm{N})\)（对每个 A-模同态 \(u: \mathrm{M} \to \mathrm{N}\)），

使得下列性质成立：

\begin{enumerate}
\def\labelenumi{(\roman{enumi})}
\item
  对每个 A-模 M，我们有 \(\mathrm{F}(1_{\mathrm{M}}) = 1_{\mathrm{F(M)}}\)。
\item
  对 \(u \in \mathrm{Hom}_{\mathrm{A}}(\mathrm{M}, \mathrm{N})\) 与 \(v \in \mathrm{Hom}_{\mathrm{A}}(\mathrm{N}, \mathrm{P})\)，我们有 \(\mathrm{F}(v \circ u) = \mathrm{F}(v) \circ \mathrm{F}(u)\)。
\item
  对 \(u, v \in \mathrm{Hom}_{\mathrm{A}}(\mathrm{M}, \mathrm{N})\)，我们有 \(\mathrm{F}(u + v) = \mathrm{F}(u) + \mathrm{F}(v)\)。
\end{enumerate}

\begin{enumerate}
\def\labelenumi{\alph{enumi})}
\item
  证明对每个 A-模 M，映射 \(u \mapsto \mathrm{F}(u)\)（从 \(\operatorname{End}_{\mathrm{A}}(\mathrm{M})\) 到 \(\operatorname{End}_{\mathrm{B}}(\mathrm{F}(\mathrm{M}))\)）是一个环同态。由此导出 \(\mathrm{F}(\mathrm{A}_{s})\) 上的一个 (B,A)-双模结构。
\item
  设 M 是一个 A-模；对 \(m \in M\)，用 \(h_{m} : A_{s} \to M\) 表示同态 \(a \mapsto am\)。证明存在唯一的 B-线性同态 \(\theta_{M}\)，从 \(\mathrm{F(A_{s})} \otimes_{\mathrm{A}} \mathrm{M}\) 到 \(\mathrm{F(M)}\)，使得 \(\theta_{\mathrm{M}}(x \otimes m) = \mathrm{F}(h_{m})(x)\)（对每个 \(x \in \mathrm{F(A_{s})}\) 与 \(m \in M\)）。我们有 \(\mathrm{F}(u) \circ \theta_{\mathrm{M}} = \theta_{\mathrm{N}} \circ (1_{\mathrm{F(A_{s})}} \otimes u)\)（对每个 A-模同态 \(u : M \to N\)）。
\end{enumerate}

\begin{enumerate}
\def\labelenumi{\arabic{enumi})}
\setcounter{enumi}{6}
\tightlist
\item
  保留前面的假设。此外，设给定了
\end{enumerate}

- 一个 A-模 \(\mathrm{G}(\mathbf{N})\)（对每个 B-模 \(\mathbf{N}\)），

- 一个 A-线性同态 \(\mathrm{G}(v):\mathrm{G}(\mathrm{N})\to \mathrm{G}(\mathrm{N}^{\prime})\)（对每个 B-模同态 \(v:\mathrm{N}\rightarrow \mathrm{N}^{\prime}\)），

- 对每个 A-模 M，一个同构 \(\alpha_{\mathrm{M}}: \mathrm{M} \to \mathrm{G}(\mathrm{F}(\mathrm{M}))\)，以及

- 一个同构 \(\beta_{\mathrm{N}}: \mathrm{N} \to \mathrm{F}(\mathrm{G}(\mathrm{N}))\)（对每个 B-模 \(\mathrm{N}\)）。

设构造 G 具有与上面 (i)–(iii) 类似的性质，并且对每个 A-模同态 \(u: M \to M'\)，我们有 \(\mathrm{G}(\mathrm{F}(u)) \circ \alpha_{\mathrm{M}} = \alpha_{\mathrm{M}'} \circ u\)，而对每个 B-模同态 \(v: N \to N'\)，我们有 \(\mathrm{F}(\mathrm{G}(v)) \circ \beta_{\mathrm{N}} = \beta_{\mathrm{N}'} \circ v\)。

\begin{enumerate}
\def\labelenumi{\alph{enumi})}
\item
  设 M 与 \(M'\) 是 A-模；证明同态 \(u \mapsto \mathrm{F}(u)\) 是从 \(\operatorname{Hom}_{\mathrm{A}}(\mathrm{M}, \mathrm{M}')\) 到 \(\operatorname{Hom}_{\mathrm{B}}(\mathrm{F}(\mathrm{M}), \mathrm{F}(\mathrm{M}'))\) 的同构。
\item
  设 \(f \in \operatorname{Hom}_{\mathrm{A}}(\mathrm{M}, \mathrm{M}')\)；则 \(\mathrm{F}(f)\) 是单射（相应地，满射）当且仅当 f 是（注意 f 是单射当且仅当对每个 A-模 E 与每个非零同态 \(g : E \to M\)，我们有 \(f \circ g \neq 0\)）。
\item
  设 \(0 \to M' \xrightarrow{i} M \xrightarrow{p} M'' \to 0\) 是 A-模的一个正合列；证明序列 \(0 \to \mathrm{F}(\mathrm{M}') \xrightarrow{\mathrm{F}(i)} \mathrm{F}(\mathrm{M}) \xrightarrow{\mathrm{F}(p)} \mathrm{F}(\mathrm{M}'') \to 0\) 是正合的。
\item
  设 \((\mathrm{M}_{\alpha})_{\alpha\in\mathrm{I}}\) 是 A-模的一个族，M 是它们的直和。构造从 \(\oplus\mathrm{F}(\mathrm{M}_{\alpha})\) 到 \(\mathrm{F}(\mathrm{M})\) 的一个典范同构。
\item
  设 M 是一个 A-模。证明 \(F(M)\) 是投射的（相应地，是生成模；相应地，是有限生成的）当且仅当 M 亦然（利用 VIII, 第 92 页的习题 9）。
\end{enumerate}

\(f\) ) 证明 (B,A)-双模 \(\mathrm{F(A_s)}\) 是可逆的，且 (A,B)-双模 \(\mathrm{G(B_s)}\) 是它的一个逆。

\(g)\) 证明对每个 A-模 M，同态 \(\theta_{\mathrm{M}}: \mathrm{F}(\mathrm{A}_s) \otimes_{\mathrm{A}} \mathrm{M} \to \mathrm{F}(\mathrm{M})\)（习题 5, \(b)\) 中定义的）是双射（利用 \(c\) ) 与 \(d)\)，通过考虑一个正合列 \(\mathrm{A}^{(\mathrm{J})} \to \mathrm{A}^{(\mathrm{I})} \to \mathrm{M} \to 0\)）。

\begin{enumerate}
\def\labelenumi{\arabic{enumi})}
\setcounter{enumi}{7}
\tightlist
\item
  设 \(k\) 是一个交换环，A 是一个 \(k\)-代数。
\end{enumerate}

\begin{enumerate}
\def\labelenumi{\alph{enumi})}
\item
  证明可逆 \((\mathrm{A},\mathrm{A})_k\)-双模（No.~8）的同构类之集，赋以乘法律 \([\mathrm{P}][\mathrm{Q}] = [\mathrm{P}\otimes_{\mathrm{A}}\mathrm{Q}]\)，是一个群；我们把它记作 \(\mathcal{P}_k(\mathrm{A})\)。
\item
  设 G 是 k-代数 A 的自同构群；对 \(g \in G\)，用 \(A_{g}\) 表示左 A-模 \(A_{s}\)，赋以由 \(x \cdot a = x g(a)\) 定义的右 A-模结构。证明映射 \(g \mapsto [A_{g}]\) 是从 G 到 \(\mathcal{P}_{k}(A)\) 的群同态，其核为内自同构子群。
\item
  设 \(\mathbf{Z}\) 是 \(\mathbf{A}\) 的中心；定义一个正合列
\end{enumerate}

\[
1 \to \mathcal {P} _ {\mathrm{Z}} (\mathrm{A}) \to \mathcal {P} _ {k} (\mathrm{A}) \to \mathrm{G}.
\]

若代数 A 是交换的，群 \(\mathcal{P}_{k}(\mathrm{A})\) 是 G 与 \(\mathcal{P}_{\mathrm{A}}(\mathrm{A})\) 的半直积。

